\section{El formato GWI: el cartucho con autoría}
\label{sec:A02GWI}


El GWI (Game \& Watch Intermedio) es el cartucho con los datos del juego y los comandos efectuados, mas los datos de intercambio: guarda las listas que produce el editor (sprites y segmentos),
las imágenes desglosadas y el HISTORIAL de la autoría. Se guarda como un ZIP
(\texttt{cartucho\_gwi.zip}), aunque se puede descomprimir y usar la carpeta generada \footnote{Esto se decide realmente en la Fase 5 Iteración 9 (ver sección \ref{sec:F5I9Usuario})} con esta disposición (el simulador lo carga
IGUAL en ZIP que descomprimido en carpeta: la forma la detecta el
programa por el contenido, no el botón):

\begin{itemize}
	\item \texttt{cartucho\_gwi.json} en la raíz: el JSON del cartucho.
	\item La imagen de segmentos, en la raíz, con el nombre 
	\texttt{img\_s}. OBLIGATORIA también en el GWI (mismo contrato que el GWB): si el ZIP no la lleva, el GWI no carga; al guardar se avisa en consola si no hay imagen de segmentos en memoria.
	\item La imagen de fondo, en la raíz, con el nombre \texttt{img\_f}. OPCIONAL, igual que en el GWB: los cartuchos sin arte de fondo no la llevan y el ZIP prescinde de ella; sin fondo, el simulador pinta la pantalla en blanco.
	\item La imagen del selector, en la raíz, con el nombre \texttt{img\_m}. OPCIONAL, igual que en el GWB: los cartuchos sin imagen para el selector no la llevan incluida.
	\item Carpeta \texttt{segmentos/} con UN PNG por cada sprite, llamado
	\texttt{img\_NNN.png} (y \texttt{img\_FNN.png} los de las piezas
	ficticias); el campo \texttt{fichero} de cada sprite referencia su
	PNG, y la resolución por nombre base hace indiferente la carpeta al
	cargar.
\end{itemize}

El json es un formato que permite ignorar estacios y saltos de línea no significativos, por lo que se genera de forma compacta para mejorar su lectura, si no el tamaño puede ser muy grande para una cómoda visualziación.

Los campos incluidos en el json son:

% Si ya declaraste estos contadores antes, omite las dos líneas \newcounter
\newcounter{fnA02GWIuno}
\newcounter{fnA02GWIdos}

\begin{table}[H]
	\centering
	\small
	\begin{tabularx}{\textwidth}{@{} l l X @{}}
		\toprule
		\textbf{Campo} & \textbf{Tipo} & \textbf{Significado} \\
		\midrule
		\texttt{tipo} & cadena & \texttt{"GWI"} \\
		\texttt{version} & número & Versión del formato (1) \\
		\texttt{nombre} & cadena & Nombre del cartucho \\
		\texttt{estilo} & cadena & Estilo de la máquina: OBLIGATORIO en el GWI (es autosuficiente y debe poder cargar sin setup previo, F5-I19) \\
		\texttt{memoria} & objeto & OPCIONAL: \texttt{\{RAM, FLASH\}} en KB, para el porcentaje de ocupación al generar el GWC (conceptos comunes y \ref{sec:F5I20}) \\
		\texttt{resolucion}\footnotemark\setcounter{fnA02GWIuno}{\value{footnote}} & objeto & \texttt{\{ancho, alto\}} en píxeles de la pantalla LCD, dentro del rango de pantallas (F5-I16). Este dato es necesario para escalar, por lo que es obligatorio incluirlo siempre.\\
		\texttt{ColorFondo} & cadena & OPCIONAL (\texttt{"\#RRGGBB"}, F5-I14): el color del fondo de la pantalla cuando el cartucho no lleva img\_f (blanco si falta); solo viaja si no es el blanco por defecto \\
		\texttt{botones} & array & \texttt{[\{numero, id, tecla?\}]} por botón (la \texttt{tecla}, F5-I17, sigue la misma regla que en el setup) \\
		\texttt{codigoGWC}\footnotemark\setcounter{fnA02GWIdos}{\value{footnote}} & cadena & OPCIONAL: El Código GWC del par, un texto con formato \texttt{"\#HHHH\#"}. Se guarda al generar el GWC y GWI a la vez; si solo se guarda el GWI, se genera sin este campo \\
		\texttt{imagenes} & objeto & \texttt{\{img\_m, img\_f, img\_s\}} (igual que el GWB) \\

		\texttt{combinar} & booleano & OPCIONAL (F3-I1): la regla «combinar» del juego (dos acciones de teclas en un mismo tick). Solo viaja si está activa; de aquí toma el GWC el bit de Reglas del Bloque 0 \\
		\texttt{fin} & objeto & OPCIONAL: \texttt{\{mecanismo, contador\}} con \texttt{mecanismo} en \texttt{vida}, \texttt{muerte} o \texttt{contador} (contador solo con el tercero, 1--4). Propiedad del cartucho, no un segmento \\
		\texttt{secuencia} & objeto & OPCIONAL: \texttt{\{ticks, segmentos\}}: animación periódica, hasta 5 segmentos lanzados en orden, uno por tick \\
		\texttt{pausas} & objeto & OPCIONAL: \texttt{\{simulador, externo\}} en ms; ausente o sin pausa: 500 ms \\
		\texttt{aceleracion} & objeto & OPCIONAL: \texttt{\{cada, resta, minimo\}}: cada \texttt{cada} ticks la pausa vigente baja \texttt{resta} ms sin pasar de \texttt{minimo} \\
		\texttt{rutas} & objeto & OPCIONAL: \texttt{\{aumentarCada\}}: cada N rutas lanzadas el máximo de rutas simultáneas sube en 1 (ausente: máximo 1) \\
		\texttt{sprites} & array & Los sprites, ver más abajo \\
		\texttt{segmentos} & array & Los segmentos, ver más abajo \\
		\texttt{autoria} & objeto & El historial de la autoría, ver más abajo \\
		\bottomrule
	\end{tabularx}
	\caption{Campos del GWI}
	\label{tab:A02GWI}
\end{table}

\footnotetext[\value{fnA02GWIuno}]{Este campo se incluyó en la Fase 5 para tener emparejados GWI y GWC.}
\footnotetext[\value{fnA02GWIdos}]{Este campo se incluyó en la Fase 5 para tener todos los datos necesarios para poder generar el GWC completo tras cargar el GWI sin necesidad de informar nada manualmente.}

Al serializar, los campos se escriben SIEMPRE en este orden (los
opcionales solo si viajan): \texttt{tipo}, \texttt{version},
\texttt{codigoGWC}?, \texttt{nombre}, \texttt{estilo},
\texttt{resolucion}, \texttt{ColorFondo}?, \texttt{memoria}?,
\texttt{imagenes}, \texttt{combinar}?, \texttt{fin}?, \texttt{secuencia}?,
\texttt{pausas}?, \texttt{aceleracion}?, \texttt{rutas}?,
\texttt{botones}, \texttt{sprites}, \texttt{segmentos},
\texttt{autoria}.

\begin{lstlisting}[
	language=json, 
	frame=single, 
	breaklines=true, 
	basicstyle=\ttfamily\small,
	caption={Ejemplo de json para un cartucho GWI (compacto y abreviado)},
	label={lst:A02GWIejemplo}
	]
	{
	"tipo": "GWI", "version": 1, "codigoGWC": \char34\#3F7E\#\char34,
	"nombre": "Parachute", "estilo": "2",
	"resolucion": { "ancho": 156, "alto": 156 },
	"memoria": { "RAM": 216, "FLASH": 128 },
	"imagenes": { "img\_m": "parachute\_m.png", "img\_f": "parachute\_f.png", "img\_s": "parachute\_s.png" },
	"combinar": true,
	"fin": { "mecanismo": "vida", "contador": null },
	"pausas": { "simulador": 500, "externo": 500 },
	"aceleracion": { "cada": 10, "resta": 10, "minimo": 100 },
	"rutas": { "aumentarCada": 3 },
	"botones": [
	{"numero": 0, "id": "IZQ", "tecla": "Q"},
	{"numero": 6, "alias": "A", "id": "Start", "tecla": "1"}
	],
	"sprites": [ { "id": "img\_001", "ancho": 16, "alto": 10, "fichero": "segmentos/img\_001.png" } ],
	"segmentos": [ { "id": "seg\_001", "x": 12, "y": 34, "sprite": "img\_001", "tipo": 1, "estado": "1", "cambiado": true, ... } ],
	"autoria": { "cambios": [ ... ], "ultimoOp": 145 }
	}
\end{lstlisting}

Viajan en este json \underline{todos los datos necesarios para diseñar y montar el
cartucho}, incluyendo el setup de la máquina para el que se ha diseñado. Los campos opcionales usan los valores
por defecto del motor cuando no viajan: el GWI no lleva sobras (los defaults pueden 
no escribirse, aunque en este momento se rellenan todos para poder verificarlos).

\subsection{Los segmentos}
\label{sec:A02Segmentos}

Cada segmento lleva los campos comunes \texttt{\{id, x, y, sprite, tipo,
	estado, cambiado\}}: identificador \texttt{seg\_NNN}, posición en
coordenadas de la imagen fuente, el sprite que lo dibuja, el tipo, y el
estado y el indicador de cambio que el motor de estados escribe en cada
tick (la lista ES el estado). Los enlaces válidos dependen del tipo:

\begin{table}[H]
	\centering
	\begin{tabularx}{\textwidth}{l l l X}
		\toprule
		Nº & Tipo & Enlaces propios & Otros campos \\
		\midrule
		0 & jugador & \makecell[tl]{\texttt{izquierda,}\\ \texttt{derecha,}\\ \texttt{arriba,}\\ \texttt{abajo}} & \texttt{primero} (jugador inicial),\newline \texttt{colision},\newline \texttt{salvadoContador},\newline \texttt{muerteContador} \\
		
		\cmidrule(lr){1-4}
		
		1 & juego & \makecell[tl]{\texttt{siguiente,}\\ \texttt{salvado,}\\ \texttt{muerte,}\\ \texttt{mostrarJunto}} & \texttt{ruta} (nº de ruta),\newline \texttt{contador} \\
		
		\cmidrule(lr){1-4}
		
		2 & inicio & como juego & El arranque de una ruta \\
		
		\cmidrule(lr){1-4}
		
		3 & vida & como juego & \texttt{contador} (orden en la cadena de vidas) \\
		
		\cmidrule(lr){1-4}
		
		4 & muerte & como juego & \texttt{contador} \\
		
		\cmidrule(lr){1-4}
		
		5 & dígito & \texttt{barra\_a}\ldots{}\texttt{barra\_g} & Las siete barras del «8» \\
		
		\cmidrule(lr){1-4}
		
		6 & contador & \texttt{digito\_1}\ldots{}\texttt{digito\_4} & \texttt{contadorInterno} (1--4),\newline \texttt{valorInicial},\newline \texttt{direccion} (ascendente/descendente),\newline \texttt{presentarCero} \\
		
		\cmidrule(lr){1-4}
		
		7 & cartucho & --- & La ficha de datos generales del cartucho: NO viaja en la lista (su mecanismo de fin es el campo \texttt{fin} del nivel superior) \\
		
		\cmidrule(lr){1-4}
		
		8 & adicional & \texttt{siguiente} & Pieza adicional \\
		
		\cmidrule(lr){1-4}
		
		9 & inicioAdicional & \texttt{siguiente} & Arranque de una cadena adicional \\
		\bottomrule
	\end{tabularx}
	\caption{Tipos de segmento y sus enlaces válidos}
	\label{tab:A02Tipos}
\end{table}

A ellos se añaden los campos transversales: \texttt{visible} (falso en los
segmentos ocultos: absorbidos y piezas internas), \texttt{fusionadoPor} (el
identificador de la pieza ficticia que lo absorbió: viaja \texttt{marcado} y el
compilador del GWC lo eliminará), y los \texttt{inversos} calculados por la
verificación al guardar: \texttt{colisionDe} (inverso de \texttt{colision})
y \texttt{mostradoJunto} (inverso de \texttt{mostrarJunto}). 

Los inversos
viajan de forma deliberadamente redundante: el compilador trabaja en el PC
sin cálculo extra. Los enlaces apuntan a identificadores legibles
(\texttt{seg\_NNN} o \texttt{FNN}); un enlace nulo (\texttt{null}) significa
sin enlace.

\subsection{Los sprites}
\label{sec:A02Sprites}

Cada sprite es \texttt{\{id, ancho, alto, fichero\}}: el identificador
legible (\texttt{img\_NNN}, o \texttt{img\_FNN} para las piezas ficticias),
sus dimensiones en píxeles de la imagen fuente y la ruta de su PNG dentro
del ZIP. Los sprites se comparten POR IDENTIDAD cuando varios segmentos
nacen de la misma componente o pieza: la deduplicación de imágenes es parte
del formato desde el origen (una imagen, una entrada en la lista, un PNG).


\subsection{La autoría: el historial}
\label{sec:A02Autoria}

La sección \texttt{autoria} guarda \texttt{\{cambios, ultimoOp\}}: el
registro CRONOLÓGICO de acciones es la única fuente de verdad de la
autoría. Cada entrada lleva al menos \texttt{\{n, relacionados\}} (número de
acción, que CONTINÚA entre guardados, y nº de elementos tocados) más los
campos de la acción con sus valores nuevos; el par
\texttt{\{segmento, campo, antes\}} describe el cambio y \texttt{antes}
conserva el valor previo (\texttt{null} si el campo nace). Al cargar un GWI,
el historial se REPRODUCE en orden sobre el estado de nacimiento (la
geometría extraída de \texttt{img\_s}): así el fichero es replay-safe y
deshacer/rehacer funciona tras recargar. Las altas y bajas de piezas
ficticias viajan como acciones (\texttt{altaFicticio}, \texttt{bajaFicticio})
con su caja y sus componentes absorbidas.

\subsection{Serialización}
\label{sec:A02Serializacion}

El JSON se escribe legible pero COMPACTO: los arrays largos (sprites,
segmentos, cambios) van a UNA línea por elemento, el resto de campos a una
línea por campo, todo en UTF-8. 

El objetivo es doble: que el fichero sea
diff-friendly (una acción de la autoría, una línea del historial) y que un
diseñador pueda de manera sencilla crearlo o editarlo a mano, si lo ve 
necesario (aunque no se recomienda, el sistema está pensado para usar el 
